
\begin{exercise}[Monte Carlo \Coffeecup]
    \label{ex: monte carlo}
    In this exercise we will devise a way to compute the integral
    \[
        I \coloneq \int_{[0,1]^\dims} f(x) dx
    \]
    of a continuous function \(f\colon [0,1]^\dims \to \real\).  Let
    \((U_n)_{n\in \nat}\) be independent and uniformly distributed on
    \([0,1]^\dims\).
    \begin{enumerate}[label=(\alph*), noitemsep]
        \item Prove that \(\hat I_n \coloneq \frac1n\sum_{i=1}^n f(U_i) \overset\as\to I\).
        \begin{solution}
            Since \(I= \E{f(U)}\) this is the strong law of large numbers,
        \end{solution}

        \item Using \(\norm{f}_\infty = \sup_{x\in [0,1]^\dims} \abs{f(x)}\) Prove that
        \[
            \prob[\big]{\abs{\hat I_n - I} > \epsilon} \le 2\exp\Bigl(-\frac{n \epsilon^2}{2\norm{f}_\infty^2}\Bigr)
        \]
        \begin{solution}
            Since \(f\) is continuous, it is bounded on the compact set \([0,1]^\dims\)
            and thereby \(f(U) \in [-\norm{f}_\infty, \norm{f}_\infty]\).
            By Hoeffding's inequality (Corollary \ref{cor: hoeffding inequality for bounded random variables})
            for bounded random variables we therefore have the claim.
        \end{solution}

        \item Let \(n = n(\epsilon, \delta)\) be the smallest \(n\) such that the following
        PAC bound holds
        \[
            \prob[\big]{\abs{\hat I_n - I} > \epsilon} \le \delta
        \]
        Prove that
        \[
            n(\epsilon, \delta)
            \le \ceil[\Big]{2\norm{f}_\infty^2\frac{\ln(\tfrac2\delta)}{\epsilon^2}}
            \in \bigO\Bigl(\norm{f}_\infty^2\frac{\log(\frac1\delta)}{\epsilon^2}\Bigr)
        \]
        \begin{remark}[Accuracy vs confidence]
            \label{rem: accuracy vs confidence}
            The accuracy \(\epsilon\) is expensive and grows like
            \(\frac1{\epsilon^2}\) while the confidence \(1-\delta\) is cheap and grows only
            like \(\log(\frac1\delta)\). Note that deterministic numerical quadrature
            formulas suffer from the curse of dimensionality causing a numerical
            complexity of order \(\epsilon^{-\frac{\dims}k}\) \citep[Sec.~1.3.12, Prop.~1]{novakDeterministicStochasticError1988}, where \(k\) is the
            number of continuous derivatives of \(f\). In comparison, the 
            computational complexity of the Monte Carlo method is independent of
            the dimension \(\dims\). This means that Monte Carlo methods become
            interesting for high-dimensional integration problems -- specifically,
            when \(\frac{\dims}k \ge 2\).
        \end{remark}
        \begin{solution}
            By the previous exercise it is sufficient to have
            \begin{align*}
                2\exp\Bigl(-\frac{n\epsilon^2}{2\norm{f}_\infty^2}\Bigr) \le \delta
                &\iff -\frac{n\epsilon^2}{2\norm{f}_\infty^2} \le \ln(\tfrac{\delta}2)
                \\
                &\iff \ln(\tfrac{2}\delta) \le \frac{n \epsilon^2}{2\norm{f}_\infty^2}
                \\
                &\iff 2\norm{f}_\infty^2 \frac{\ln(\tfrac2\delta)}{\epsilon^2} \le n
            \end{align*}
            Since this is a sufficient condition, the smallest \(n(\delta, \epsilon)\)
            is smaller than the first natural number that satisfies it, leading
            to the claim.
        \end{solution}
    \end{enumerate}
\end{exercise}
